\documentclass[11pt]{article}
\usepackage[margin=0.85in]{geometry}
\usepackage{amsmath,amssymb,mathtools,booktabs,microtype,xcolor,fancyhdr,hyperref}
\definecolor{navy}{HTML}{17365D}\definecolor{teal}{HTML}{147D78}
\hypersetup{colorlinks=true,linkcolor=navy,urlcolor=teal}
\pagestyle{fancy}\setlength{\headheight}{14pt}\fancyhf{}\lhead{Qudit Pauli and finite Heisenberg groups}\rhead{Proof and verification}\cfoot{\thepage}
\newcommand{\Sgm}{\mathcal S}
\newcommand{\E}{\mathbb E}
\newcommand{\F}{\mathbb F}
\title{\color{navy}\bfseries Qudit Pauli and Finite-Heisenberg Sigma--MGFs\\\large Proof, matrix construction, and direct checks}
\author{Computational verification note}\date{17 July 2026}
\begin{document}\maketitle

\begin{abstract}
For odd characteristic we prove the two-spectrum formula for the defining
Schrödinger representation and check it by constructing every matrix in
$P_{3,1}$, $P_{5,1}$, $P_{9,1}$, and $H_3(\F_9)$.  The tests evaluate the
target invariant dimension directly, not from the spectral classification
used in the proof.  All 17 representative coefficient comparisons pass.
\end{abstract}

\section{Statement}
Let $q=p^f$ with $p$ odd, $V=\mathbb C[\F_q^n]$, and $N=q^n$.  Fix a
nontrivial additive character $\psi:\F_q\to\mu_p$.  With
\[
X(a)e_x=e_{x+a},\qquad Z(b)e_x=\psi(b\cdot x)e_x,
\qquad W(a,b)=\psi\!\left(\tfrac12a\cdot b\right)X(a)Z(b),
\]
put $P_{q,n}=\{\zeta W(a,b):\zeta\in\mu_p\}$.  For a partition
$\tau=(\tau_1,\ldots,\tau_\ell)$ define
\[
A_\tau=\prod_i\binom{N+\tau_i-1}{\tau_i},\qquad
B_\tau=\begin{cases}
\displaystyle\prod_i\binom{N/p+\tau_i/p-1}{\tau_i/p},&p\mid\tau_i\ \forall i,\\
0,&\text{otherwise.}
\end{cases}
\]
Then the claimed generalized Molien coefficient is
\begin{equation}\label{eq:coefficient}
[m_\tau]\Sgm_{P_{q,n}}=
\begin{cases}
0,&p\nmid |\tau|,\\
q^{-2n}A_\tau+(1-q^{-2n})B_\tau,&p\mid |\tau|.
\end{cases}
\end{equation}
Equivalently,
\begin{equation}\label{eq:series}
\Sgm_{P_{q,n}}(\mathbf t)=
\frac{q^{-2n}}p\sum_{\zeta\in\mu_p}\prod_j(1-\zeta t_j)^{-N}
+(1-q^{-2n})\prod_j(1-t_j^p)^{-N/p}.
\end{equation}

\section{Proof}
There are $p$ central matrices $\zeta I_N$, a fraction $q^{-2n}$ of the
group.  Their average is the first term of \eqref{eq:series}.  Now take
$(a,b)\ne(0,0)$.  For $1\le k<p$, $W(a,b)^k$ has trace zero: if $ka\ne0$
it is a translation with zero diagonal; otherwise its diagonal trace is a
nontrivial additive-character sum.  Also $W(a,b)^p=I$.  Fourier inversion of
the multiplicities of the $p$th roots of unity therefore gives multiplicity
$N/p$ for every root.  Hence
\[
\det(I-tW(a,b))=(1-t^p)^{N/p}.
\]
A scalar in $\mu_p$ only permutes this spectrum, proving \eqref{eq:series}.
Expanding the central factors gives $A_\tau$ and forces $p\mid|\tau|$;
expanding the noncentral factors gives $B_\tau$.  This proves
\eqref{eq:coefficient}.

For the finite Heisenberg group $H_{2n+1}(\F_q)$, the Schrödinger matrix is
$\rho_\psi(a,b,z)=\psi(z)W(a,b)$.  When $f>1$, $\ker\psi$ has size $q/p$,
so each matrix in $P_{q,n}$ has exactly $q/p$ preimages.  Uniform averages
over the abstract group and its image consequently agree.

\section{Independent verification method}
For every constructed matrix $g$, the code computes its eigenvalues
$\lambda_1,\ldots,\lambda_N$ numerically and obtains $h_d(\lambda)$ from
\[
\prod_{r=1}^N(1-\lambda_r t)^{-1}=\sum_{d\ge0}h_d(\lambda)t^d.
\]
It then evaluates the target directly as
\[
D_\tau=\frac1{|G|}\sum_{g\in G}\prod_i h_{\tau_i}(\operatorname{spec}g).
\]
This is the character of $\bigotimes_i\operatorname{Sym}^{\tau_i}V$ averaged
over $G$, hence the rank of its Reynolds projector.  Only after computing
$D_\tau$ does the program evaluate \eqref{eq:coefficient}.

The $q=9$ construction uses $\F_9=\F_3[u]/(u^2+1)$ and
$\psi(x)=\exp(2\pi i\operatorname{Tr}_{\F_9/\F_3}(x)/3)$.  Thus the test
includes a genuine prime-power field.  The $H_3(\F_9)$ run enumerates all
729 abstract triples, retaining the threefold central kernel rather than
silently replacing the group by its image.

\section{Representative results}
\begin{center}
\begin{tabular}{@{}lrrrr@{}}\toprule
group & $\dim V$ & matrices averaged & $\tau$ & direct $=$ formula\\\midrule
$P_{3,1}$ &3&27&$(3)$&2\\
$P_{3,1}$ &3&27&$(3,3)$&12\\
$P_{5,1}$ &5&125&$(4,1)$&14\\
$P_{5,1}$ &5&125&$(5,5)$&636\\
$P_{9,1}$ &9&243&$(3)$&5\\
$P_{9,1}$ &9&243&$(3,3)$&345\\
$H_3(\F_9)$ &9&729&$(2,1)$&5\\
$H_3(\F_9)$ &9&729&$(3,3)$&345\\\bottomrule
\end{tabular}
\end{center}
All 17 built-in cases pass to tolerance $2\times10^{-8}$ and round to
nonnegative integers.  Characteristic two is intentionally excluded here:
the displayed half-phase normalization is an odd-characteristic statement;
the phase-extended qubit lift is handled in its own verification.

\section{Reproduction}
Run
\begin{verbatim}
.venv/bin/python lambda_distributions/proofs/quantum_matrix_groups/
  pauli_heisenberg/verification.py
.venv/bin/marimo edit marimo_notebooks/pauli_heisenberg.py
\end{verbatim}
The line breaks after the directory slash are for display; join each command
onto one shell line.
\end{document}
